% This work is licensed under the Creative Commons
% Attribution-NonCommercial-ShareAlike 4.0 International License. To view a copy
% of this license, visit http://creativecommons.org/licenses/by-nc-sa/4.0/ or
% send a letter to Creative Commons, PO Box 1866, Mountain View, CA 94042, USA.
% (c) Eric Kunze, 2020
%%%%%%%%%%%%%%%%%%%%%%%%%%%%%%%%%%%%%%%%%%%%%%%%%%%%%%%%%%%%%%%%%%%%%%%%%%%%%%%
% Template for exercises and homework at TU Dresden.
%%%%%%%%%%%%%%%%%%%%%%%%%%%%%%%%%%%%%%%%%%%%%%%%%%%%%%%%%%%%%%%%%%%%%%%%%%%%%%%

\documentclass[ngerman, a4paper, 11pt]{article}

\usepackage[ngerman]{babel}
\usepackage[top=2.5cm,bottom=2.5cm,left=2.5cm,right=2.5cm]{geometry}

\usepackage{parskip}    % split paragraphs by vspace instead of intendations
\usepackage[onehalfspacing]{setspace} % increase row-space

\usepackage{xifthen}
\usepackage{xparse}

%%%%%%%%%%%%%%%%%%%%%%%%%%%%%%%%%%%%%%%%%%%%%%%%%%%%%%%%%%%%%%%%%%%%%%%%%%%%%%%
% SCHRIFTEN
\usepackage[utf8]{inputenc}
\usepackage{chngcntr}
\usepackage{eufrak}

\usepackage{lmodern}
\usepackage[normalem]{ulem}

\usepackage[autostyle=true,english=british]{csquotes}

% new font OpenSans
\usepackage[scale=1]{opensans}
\newcommand*{\osfamily}{\fontfamily{fos}\selectfont}
\DeclareTextFontCommand{\textos}{\osfamily}

\usepackage{listings}

%%%%%%%%%%%%%%%%%%%%%%%%%%%%%%%%%%%%%%%%%%%%%%%%%%%%%%%%%%%%%%%%%%%%%%%%%%%%%%%
% EIGENSCHAFTEN
\newcommand{\name}{Eric Kunze}
\newcommand{\matnr}{4679202}
\newcommand{\email}{\href{mailto:eric.kunze@mailbox.tu-dresden.de}{\ttfamily eric.kunze@mailbox.tu-dresden.de}}

\newcommand{\modul}{Wahrscheinlichkeitstheorie mit Martingalen}
\newcommand{\semester}{Wintersemester 2020/21}

%\renewcommand{\tutor}{TUTOR}
%\renewcommand{\gruppe}{Tag x. DS, (un)gerade Woche}

\newcommand{\professor}{Prof. Dr. Martin Keller-Ressel}
\newcommand{\fakultaet}{Mathematik}
\newcommand{\institut}{Stochastik}
\newcommand{\lehrstuhl}{Stoch. Analysis \& Finanzmathematik}

%%%%%%%%%%%%%%%%%%%%%%%%%%%%%%%%%%%%%%%%%%%%%%%%%%%%%%%%%%%%%%%%%%%%%%%%%%%%%%%
% GRAPHICS
\usepackage{graphicx}
\usepackage[table,dvipsnames]{tudscrcolor} % package xcolor loaded automatically
\usepackage[font=small,labelfont=bf]{caption} % captions of non-floated figures

\usepackage{tcolorbox}

\usepackage{tikz}
\usetikzlibrary{matrix, patterns,arrows,calc,decorations.pathmorphing,backgrounds, positioning,fit,petri,decorations.fractals}
\usepackage{pgf}
\usepackage{pgfplots}
\pgfplotsset{compat=1.10}   % in my packages used compat=1.15
\usepgfplotslibrary{fillbetween}

% tabulars
\usepackage{tabularx}   % tabularx-environment (explicitly set width of columns)
\usepackage{multirow}
\usepackage{booktabs}   % improved rules

%%%%%%%%%%%%%%%%%%%%%%%%%%%%%%%%%%%%%%%%%%%%%%%%%%%%%%%%%%%%%%%%%%%%%%%%%%%%%%%
% ENUMERATIONS
\usepackage{enumerate}
\usepackage[inline]{enumitem}       % customize label

\renewcommand{\labelitemi}{\raisebox{1pt}{\scalebox{.4}{$\blacksquare$}}}
\renewcommand{\labelitemii}{$\vartriangleright$}
\renewcommand{\labelitemiii}{--}
% Variantionen des Dreiecks als Aufzählungszeichen $\blacktriangleright$ / $\vartriangleright$ / $\triangleright$

\renewcommand{\labelenumi}{(\arabic{enumi})}
\renewcommand{\labelenumii}{\alph{enumii}.}
\renewcommand{\labelenumiii}{\roman{enumiii}.}

%%%%%%%%%%%%%%%%%%%%%%%%%%%%%%%%%%%%%%%%%%%%%%%%%%%%%%%%%%%%%%%%%%%%%%%%%%%%%%%
% PAGE LAYOUT
\usepackage{scrlayer-scrpage}
\clearpairofpagestyles
\newpairofpagestyles[scrheadings]{firstpage}{\ohead*{}\cfoot*{\pagemark}}
\newpairofpagestyles[scrheadings]{normalpage}{
    \ihead*{\fosfamily\color{cdgray}\small\modul}
    \ohead*{\fosfamily\color{cdgray}\small\name~ (\matnr)}
    \cfoot*{\pagemark}}

\pagestyle{normalpage}
\pagenumbering{arabic}

%%%%%%%%%%%%%%%%%%%%%%%%%%%%%%%%%%%%%%%%%%%%%%%%%%%%%%%%%%%%%%%%%%%%%%%%%%%%%%%
% TITLEPAGE
\newcommand{\makeTUtitle}[1][]{%
    \begin{titlepage}
        \pagecolor{cddarkblue!90} \color{white}%
        \raggedright \fosfamily%
        \setlength{\parindent}{0pt}%
    % Logo / Kopf
        \hspace{-18.6mm} %
        \includegraphics[scale=0.6]{TUD-white.pdf} \\
        \vspace{3mm}
        \begin{tabular}{m{\textwidth}}
            \hline
            \hspace{-4pt}\small{\textbf{Fakultät \fakultaet} Institut für \institut, Professur für \lehrstuhl} \\
            \hline
        \end{tabular} \\
    % Titel
        \vspace{5cm}
        {\Huge\bfseries \MakeUppercase \modul \par}
        \vspace{0.5cm}%
        {\Large \itshape Übungen} \\%
        \vspace{1.5cm}
        \textbf{{\Large \professor}} \par
        \vspace{0.5cm}
        {\large \semester}
    % Fußzeile
        \vfill%
        \begin{tabular}{lll}
            Autor  & : & \name \\
            E-Mail & : & \email \\
        \end{tabular}%
    \end{titlepage}
    \nopagecolor
}

%%%%%%%%%%%%%%%%%%%%%%%%%%%%%%%%%%%%%%%%%%%%%%%%%%%%%%%%%%%%%%%%%%%%%%%%%%%%%%%
% COUNTER
\usepackage{chngcntr}

\newcounter{taskcount}
\newcounter{blattcount}
\newcounter{thmcount}

\counterwithin{page}{blattcount}
\counterwithin{equation}{blattcount}
\counterwithin{figure}{blattcount}
\counterwithin{table}{blattcount}

% automatic reset of task counter in each blatt
%\pretocmd{\exercisePage}{\setcounter{taskcount}{0}}{}{}

%\counterwithin{taskcount}{blattcount}

%%%%%%%%%%%%%%%%%%%%%%%%%%%%%%%%%%%%%%%%%%%%%%%%%%%%%%%%%%%%%%%%%%%%%%%%%%%%%%%
% EXERCISE PAGE
\newcommand{\header}{%
        {\fosfamily%
            \fcolorbox{cddarkblue}{cdblue!20}{%
                \begin{minipage}{\dimexpr0.74\linewidth-\fboxrule-\fboxsep}
                    {\huge \textbf{Hausaufgaben}} \vspace{3pt} \\
                    \textbf{\modul} -- Übungsblatt \theblattcount
                \end{minipage}%
                \begin{minipage}{\dimexpr0.24\linewidth-\fboxrule-\fboxsep}
                    \flushright \textbf{\name} \\
                    Matr.-Nr. \matnr
                \end{minipage}%
        }}%
    }%

\NewDocumentEnvironment{exercisePage}{O{}}{%
    \pagebreak%
    \stepcounter{blattcount} \setcounter{page}{1}%
    \thispagestyle{firstpage}
    \header \\[6pt]
    \setcounter{figure}{0}
    \setcounter{equation}{0}%
    \setcounter{table}{0}%
    \setcounter{thmcount}{0}
    %
    \ifthenelse{\isempty{#1}}{}{%
        {\fosfamily \itshape Thema: #1} \par%\\[12pt]%
    }
}{}

%%%%%%%%%%%%%%%%%%%%%%%%%%%%%%%%%%%%%%%%%%%%%%%%%%%%%%%%%%%%%%%%%%%%%%%%%%%%%%%
% MATH ENVIRONMENTS
\usepackage{../mathoperatorsMathTUD}

\usepackage[ntheorem,framemethod=TikZ]{mdframed}
\usepackage{amsmath,amssymb,amsfonts,mathtools}
\usepackage[amsmath,thmmarks,framed]{ntheorem}

\setlength\abovedisplayshortskip{0pt plus 3pt}%
\setlength\belowdisplayshortskip{4pt plus 2pt minus 2pt}%
\setlength\abovedisplayskip{1pt plus 1pt minus 2pt}%
\setlength\belowdisplayskip{1pt plus 1pt minus 2pt}%

\newcommand{\skiparound}{10pt}

% >> Exercises
\theoremstyle{plain}
\theoremheaderfont{\fosfamily\normalsize\bfseries\upshape}
\theorembodyfont{\normalsize}
\theoremseparator{.}
\theoremsymbol{}

\newmdtheoremenv[%
    backgroundcolor=cdblue!5,%
    linecolor=cddarkblue,%
    skipabove=\skiparound,%
    skipbelow=\skiparound,%
    nobreak,%
]{exercise}[taskcount]{Übung}

\newmdtheoremenv[%
    backgroundcolor=cdblue!5,%
    linecolor=cddarkblue,%
    skipabove=\skiparound,%
    skipbelow=\skiparound,%
    nobreak,%
]{homework}[taskcount]{Hausaufgabe}

\newmdtheoremenv[%
    backgroundcolor=cdblue!5,%
    linecolor=cddarkblue,%
    skipabove=\skiparound,%
    skipbelow=\skiparound,%
    nobreak,%
]{task}[taskcount]{Aufgabe}

\newmdtheoremenv[%
backgroundcolor=cdblue!5,%
linecolor=cddarkblue,%
skipabove=\skiparound,%
skipbelow=\skiparound,%
nobreak,%
]{zusatz}[taskcount]{Zusatzaufgabe}

% >> Lemma
\theoremstyle{plain}
\theoremheaderfont{\fosfamily\normalsize\bfseries\upshape}
\theorembodyfont{}
\theoremseparator{.}
\theorempreskip{6pt}
\theorempostskip{6pt}
\newtheorem{lemma}{Lemma}
\counterwithin{lemma}{blattcount}

% >> proof
\makeatletter
\newtheoremstyle{proofstyle}%
{\item[\hskip\labelsep {\theorem@headerfont ##1}\theorem@separator]}%
{\item[\hskip\labelsep {\theorem@headerfont ##1}\ (##3)\theorem@separator]}
\makeatother

\theoremstyle{proofstyle}
\theoremheaderfont{\normalsize\slshape}
\theorembodyfont{}
\theoremseparator{.}
\theorempreskip{0pt}
\theorempostskip{5pt}
\theoremsymbol{$\square$}
\newtheorem{proof}{Beweis}

% >> solutions
\newtheorem{solution}{Lösung}

% >> Equivalences
\newcommand{\hinrichtung}{\item[\bfseries ($\boldsymbol{\Rightarrow}$)]}
\newcommand{\rueckrichtung}{\item[\bfseries ($\boldsymbol{\Leftarrow}$)]}
\newenvironment{equivalence}[1][]{%
        \ifthenelse{\isempty{#1}}%
        {\begin{description}} % no optional argument
            {\begin{description}[topsep=-\parskip]} % any optional argument
    }{ \end{description} }%

    % >> inductions (needs description environment)
    \newcommand{\ianfang}[1][]{%
        \ifthenelse{\isempty{#1}}{%
            \item[\textbf{(IA)}] % no parameter
        }{  \item[\textbf{(IA)}] {#1 :}  } % parameter exists
    }

    \newcommand{\ivorraussetzung}{\item[\bfseries (IV)]}

    \newcommand{\ischritt}[1][]{%
        \ifthenelse{\isempty{#1}}{%

            \item[\textbf{(IS)}] % no parameter
        }{  \item[\textbf{(IS)}] {#1 :} } % parameter exists
    }

%%%%%%%%%%%%%%%%%%%%%%%%%%%%%%%%%%%%%%%%%%%%%%%%%%%%%%%%%%%%%%%%%%%%%%%%%%%%%%%
% HYPERLINKS
\usepackage[unicode,bookmarks=true]{hyperref}
\hypersetup{
    % pdfborder={0 0 0}         % no boxed around links
    pdfborderstyle={/S/U/W 1},  % underlining insteas of boxes
    linkbordercolor=cdblue,
    urlbordercolor=cdblue
    %   colorlinks,
    %   citecolor=black,
    %   filecolor=cddarkblue!80,
    %   linkcolor=black,
    %   urlcolor=cddarkblue!80
}
\usepackage{bookmark}       % pdf-bookmarks

\usepackage{cleveref}
\crefname{beispiel}{Beispiel}{Beispiele}
\crefname{erinnerung}{Erinnerung}{Erinnerungen}
\crefname{wiederholung}{Wiederholung}{Wiederholungen}
\crefname{bemerkung}{Bemerkung}{Bemerkungen}
\crefname{enumi}{Teil}{Teile}

%%%%%%%%%%%%%%%%%%%%%%%%%%%%%%%%%%%%%%%%%%%%%%%%%%%%%%%%%%%%%%%%%%%%%%%%%%%%%%%
% HIGHLIGHTING
\newcommand{\begriff}[1]{\textbf{#1}}
\newcommand{\person}[1]{\textsc{#1}}

%%%%%%%%%%%%%%%%%%%%%%%%%%%%%%%%%%%%%%%%%%%%%%%%%%%%%%%%%%%%%%%%%%%%%%%%%%%%%%%
% ADDITIONAL COMMANDS
\renewcommand{\Var}{\operatorname{\mathbb Var}}
\renewcommand{\Cov}{\operatorname{\mathbb Cov}}

\renewcommand{\F}{\mathcal{F}}

%%%%%%%%%%%%%%%%%%%%%%%%%%%%%%%%%%%%%%%%%%%%%%%%%%%%%%%%%%%%%%%%%%%%%%%%%%%%%%%

\begin{document}

\makeTUtitle[Übungen]

\begin{exercisePage}
	
% AUFGABE 1	
	\begin{task}
		Seien $X,Y \in L^2(\mathcal{A})$ und $\F$ eine Unter-$\sigma$-Algebra von $\mathcal{A}$. Bedingte
		Varianz und bedingte Kovarianz von $X$ bzw. $X,Y$ bezüglich $\F$ sind definiert als
		\begin{align*}
			\Var(X | \F) &\defeq \EW[{(X - \EW[X | \F])^2 \mid \F}] \\
			\Cov(X,Y | \F) &\defeq \EW[ {(X  \EW[X | \F]) (Y - \EW[Y | \F]) \mid \F} ]
		\end{align*}
		Zeige die Sätze von der totalen Varianz bzw. totalen Kovarianz:
		\begin{align*}
			\Var(X) = \EW[\Var(X | \F)] + \Var(\EW[X | \F]) \\
			\Cov(X,Y) = \EW[ \Cov(X,Y | \F) ] + \Cov( \EW[X | \F] , \EW[Y | \F])
		\end{align*}
	\end{task}

	Sei $U \defeq \EW[X | \F]$ und $V \defeq \EW[Y | \F]$. Mit der Turmregel gilt $\EW[U] = \EW[X]$ bzw. $\EW[V] = \EW[Y]$. Mit dem Verschiebungssatz gilt $\Var(U) = \EW[U^2] - \EW[X]^2$ und analog für $V$.
	Offenbar sind $U$ und $V$ nun $\F$-messbar.
	Es ist
	\begin{equation*}
		\Var(X | \F) = \EW[(X - U) | \F]
		= \EW[X^2 - 2XU + U^2 | \F] 
		= \EW[X^2 | \F] - 2U \underbrace{\EW[X | \F]}_{= U} + \EW[U^2]
	\end{equation*}
	und somit
	\begin{align*}
		\EW[\Var(X | \F)] 
		= \EW[ {\EW[X^2 | \F] - 2U \EW[X | \F] + \EW[U^2]} ] 
		&= \EW[X^2] - 2 \EW[U^2] + \EW[U^2] \\
		&= \EW[X^2] - \EW[U^2]
	\end{align*}
	Weiterhin gilt
	\begin{align*}
		\Var(U) &= \EW[ {(U- \EW[U])^2} ] = \EW[ {U^2 - 2 U \EW[X] + \EW[X]^2} ] 
		= \EW[U^2] - 2 \EW[X] \EW[U] + \EW[X]^2 \\
		&= \EW[U^2]- \EW[X]^2 
	\end{align*}
	Damit gilt schließlich 
	\begin{align*}
		\EW[\Var(X | \F)] + \Var(\EW[X | \F]) 
		&= \EW[X^2] - \EW[U^2] + \EW[U^2]- \EW[X]^2  
		= \EW[X^2] - \EW[X]^2 \\
		&= \Var(X)
	\end{align*}
	
	Analog verfahren wir beim Satz von der totalen Kovarianz:
	Es gilt
	\begin{align*}
		\EW[ \Cov(X,Y | \F) ] 
		&= \EW[ {\EW[ (X-U)(Y-V) \mid \F]} ] \\
		&= \EW[ {\EW[ XY | \F] - \EW[VX | \F] - \E[UY | \F] + \EW[UV | \F]} ] \\
		&= \EW[XY]- \E[UV]- \E[UV]+ \E[UV] \\
		&= \E[XY]- \E[UV]
	\intertext{sowie}
		\Cov(U,V) 
		&= \EW[{(U-\EW[U])(V-\EW[V])}] \\
		&= \EW[{UV- U*\EW[Y] - V*\EW[X] + \EW[X]*\EW[Y]}] \\
		&= \EW[UV] - \EW[X]*\EW[Y] - \EW[X]*\EW[Y] + \EW[X]*\EW[Y]  \\
		&= \EW[UV] - \EW[X]*\EW[Y]
	\end{align*}
	und damit schlussendlich
	\begin{equation*}
		\Cov\left( \EW[X | \F], \EW[Y | \F] \right) + \EW[ \Cov(X,Y | \F)] = \EW[XY] - \EW[X]\EW[Y] = \Cov(X,Y)
	\end{equation*}
	
	
% AUFGABE 2	
	\begin{task}
		\begin{enumerate}[label=(\alph*)]
			\item In einer bestimmten Population sei das Alter $X$ bei erstmaliger Berufsunfähigkeit exponentialverteilt mit Parameter $\lambda > 0$. Für eine Versicherungsgesellschaft die gegen Berufsunfähigkeit versichert, ist das mittlere Alter bei Eintritt der Berufsunfähigkeit von Bedeutung, unter der Bedingung, dass die Berufsunfähigkeit zwischen den Altersgrenzen $0 \le a < b$ eintritt. Bestimme diesen bedingten Erwartungswert $\EW[X | a \le X \le b]$.
			\item Welche der folgenden in der Vorlesung definierten mathematischen Objekte sind reelle Zahlen, Zufallsvariablen oder messbare Funktionen von $\R$ nach $\R$?
			\begin{equation*}
				\EW[X|\F], \quad \P(A|B), \quad \EW[X|Y=y], \quad \P(A|\F), \quad \EW[X|Y]
			\end{equation*}
			Wie üblich bezeichnet $\F$ eine $\sigma$-Algebra, $X$ und $Y$ sind reellwertige Zufallsvariablen, $y$ eine reelle Zahl und $A$ und $B$ sind Ereignisse.
		\end{enumerate}
	\end{task}

	\begin{enumerate}[label=(\alph*)]
		\item Es sei $X \sim \Exp(\lambda)$, d.h. $X$ hat die Dichte $f(x) = \lambda e^{-\lambda x} * \one_{[0,\infty)}(x)$.
		Da $[a,b] \subseteq [0,\infty]$ können wir rechnen
		\begin{align*}
			\EW[X \mid a \le X \le b] &= \int_\R x * \lambda e^{-\lambda x} *\one_{[0,\infty)}(x) * \one_{[a,b]}(x) \dx 
			= \int_a^b x * \lambda e^{-\lambda x} \dx \\
			&= \sqbrackets{-x * e^{-\lambda x}}_a^b + \int_a^b e^{-\lambda x} \dx \tag{partielle Integration} \\
			&= \sqbrackets{-x * e^{-\lambda x}}_a^b - \sqbrackets{\frac{1}{\lambda} e^{-\lambda x}}_a^b \\
			&= \brackets{a + \frac{1}{\lambda}} * e^{-\lambda a} - \brackets{b + \frac{1}{\lambda}} * e^{-\lambda b}
		\end{align*}
		\item Es sind
		\begin{itemize}
			\item $\EW[X | \F]$ und $\EW[X | Y]$ Zufallsvariablen,
			\item $\P(A | B)$ und $\P(A | \F)$ reelle Zahlen sowie
			\item $\EW[X | Y=y]$ eine messbare Funktion von $\R$ nach $\R$.
		\end{itemize}
	\end{enumerate}


% AUFGABE 3
	\newcommand{\xy}{\left(\begin{smallmatrix} x \\ y \end{smallmatrix}\right)}
	
	\begin{task}
		Sei $M = \left(\begin{smallmatrix} A & B \\ C & D \end{smallmatrix} \right)$ eine quadratische, in Blöcke unterteilte Matrix von vollem Rang, wobei $A$ und $D$ ebenfalls quadratisch und von vollem Rang seien. Die Ausdrücke
		\begin{equation*}
			(M/A) \defeq (D - CA^{-1} B) \quad \und \quad (M/D) \defeq (A - BD^{-1} C)
		\end{equation*}
		heißen \textit{Schurkomplement} von $A$ in $M$ bzw. $D$ in $M$. Folgende Formel für die blockweise Inversion von $M$ dürfen Sie als gegeben betrachten:
		\begin{equation*}
			M^{-1} = \begin{pmatrix}
				(M/D)^{-1} & -A^{-1}B(M/A)^{-1} \\
				-D^{-1}C(M/D)^{-1} & (M/A)^{-1}
			\end{pmatrix}
		\end{equation*}
		Verwenden Sie die Formel um folgende Teilaufgaben zu lösen:
		\begin{enumerate}[label=(\alph*)]
			\item \label{it: ue1-aufgabe3a}
			Sei $M$ nun symmetrisch, d.h. $A$ und $D$ sind symmetrisch und $C = B^\top$. Zeige
			\begin{equation*}
				\brackets{x^\top, y^\top} M^{-1} \xy - y^\top D^{-1} y 
				=
				\schlange{x}^\top (M/D)^{-1} \schlange{x}
			\end{equation*}
			mit $\schlange{x} = (x - BD^{-1}y)$ für alle $x,y$ mit passender Dimension.
			\item Es sei $(X,Y)$ multivariat normalverteilt mit Erwartungswert $0$ und positiv definiter Kovarianzmatrix $\Sigma = \left(\begin{smallmatrix} \Sigma_X & \Sigma_{XY} \\ \Sigma_{XY} & \Sigma_Y \end{smallmatrix}\right)$. Dann ist $X$ bedingt auf $Y$ normalverteilt mit $\EW[X|Y] = \Sigma_{XY}\Sigma_Y^{-1}$ und Kovarianzmatrix $(\Sigma/\Sigma_Y)$. \\
			Hinweis: Es gilt $\det(\Sigma) = \det(\Sigma_Y) * \det(\Sigma/\Sigma_Y)$.
		\end{enumerate}
	\end{task}

	\begin{enumerate}[label=(\alph*)]
		\item Wir notieren $M^{-1} = T = (T_{ij})_{i,j}$. Dabei ist $T$ wieder symmetrisch, da zum Einen
		\begin{equation*}
			(M/D)^\top
			= (A - BD^{-1}C)^\top 
			= A^\top - C^\top \transpose{D^{-1}} B^\top
			= A - BD^{-1}C 
			= (M/D)
		\end{equation*}
		sowie $(M/A)^\top = (M/A)$ analog und zum Anderen
		\begin{equation}
			\begin{alignedat}{2}
				&& A^{-1} B(M/A)^{-1} &= (M/D)^{-1} BD^{-1} \\
				\Leftrightarrow &\quad& (M/D)A^{-1}B &= BD^{-1} (M/A) \\
				\Leftrightarrow &\quad& (A-BD^{-1}C)A^{-1}B  &= BD^{-1}(D-CA^{-1}B) \\
				\Leftrightarrow &\quad& B-BD^{-1}CA^{-1}B &= B-BD^{-1}CA^{-1}B 
			\end{alignedat}
			\tag{$\star$} \label{eq: ue1-aufgabe3b-taut1}
		\end{equation}
		tautologisch ist.
		Betrachten wir nun $q(x,y) = \brackets{x^\top, y^\top} T \xy - y^\top D^{-1} y$, so erhalten wir
		\begin{align*}
			q(x,y) 
			&= x^\top T_{11} x + x^\top T_{12} y + y^\top T_{21} x + y^\top (T_{22} - D^{-1}) y
		\end{align*}
		Mit einer Matrix $\Lambda$ können wir dies in die Form
		\begin{equation*}
			q(x,y) 
			= (x^\top - y^\top \Lambda^\top) T_{11} (x + \Lambda y)
			= x^\top T_{11} x - x^\top \underbrace{(T_{11} \Lambda)}_{=T_{12}} y - y^\top \underbrace{(\Lambda^\top T_{11})}_{= T_{21}} x + y^\top \underbrace{(\Lambda^\top T_{11} \Lambda)}_{= T_{22} - D^{-1}} y
		\end{equation*}
		Aus $-T_{11} \Lambda = T_{12}$ erhalten wir $\Lambda = -T_{11} T_{12} = (M/D) * A^{-1} B (M/A)^{-1} \overset{{\eqref{eq: ue1-aufgabe3b-taut1}}}{=} BD^{-1}$. Außerdem gilt damit $T_{12} = T_{21}$ und
		\begin{alignat*}{2}
			&& (M/A)^{-1}-D^{-1} &= D^{-1} C(M/D)^{-1} BD^{-1} \\
			\Leftrightarrow &\quad& (M/A)^{-1} DB^{-1} - B^{-1} &= D^{-1} C(M/D)^{-1} \\
			\Leftrightarrow &\quad& DB^{-1} (M/D) -(M/A)B^{-1} (M/D) &= (M/A)D^{-1}C \\
			\Leftrightarrow &\quad& DB^{-1}A-C-DB^{-1}A+C+C-CA^{-1}BD^{-1}C &= C- CA^{-1}BD^{-1}C
		\end{alignat*}
		tautologisch. Somit ist mit $\schlange{x} \defeq (x - BD^{-1}y)$ schließlich 
		\begin{equation*}
			q(x,y) = \schlange{x}^\top (M/D)^{-1} \schlange{x}
		\end{equation*}
		und die Behauptung bewiesen.
		
		
		
		\item Sei $f_{XY}$ die gemeinsame Dichte von $(X,Y)$ gegeben durch
		\begin{align*}
			f_{XY}(x,y) 
			&= \frac{1}{\sqrt{(2\pi)^p \det(\Sigma)}} * \exp\brackets{-\frac{1}{2} \brackets{(x^\top, y^\top) - (\mu_X^\top, \mu_Y^\top)} * \Sigma^{-1} * \brackets{\left(\begin{smallmatrix} x \\ y \end{smallmatrix}\right) - \left(\begin{smallmatrix} \mu_X \\ \mu_Y \end{smallmatrix} \right)}} \\
			&= \frac{1}{\sqrt{(2\pi)^p \det(\Sigma)}} * \exp\brackets{-\frac{1}{2} (x^\top, y^\top) * \Sigma^{-1} * \left(\begin{smallmatrix} x \\ y \end{smallmatrix}\right)}
		\end{align*}
		Die Randverteilung von $Y$ ist gegeben durch
		\begin{equation*}
			f_Y(y) = \frac{1}{\sqrt{(2\pi)^{p_Y} \det(\Sigma_Y)}} * \exp\brackets{-\frac{1}{2} y^\top * \Sigma_Y^{-1} * y}
		\end{equation*}
		Damit erhalten wir die bedingte Dichte $f_{X|Y}$ als
		\begin{align*}
			f_{X|Y=y}(x) &= \frac{f_{XY}(x,y)}{f_Y(y)} 
			= \sqrt{(2\pi)^{p - p_Y} * \frac{\det(\Sigma_Y)}{\det(\Sigma)}} * \exp\brackets{-\frac{1}{2} (x^\top, y^\top) * \Sigma^{-1} * \xy - y^\top * \Sigma_Y^{-1} * y} \\
			&= \frac{1}{\sqrt{(2\pi)^{p - p_Y} \det(\Sigma / \Sigma_Y)}} * \exp\brackets{-\frac{1}{2} (x^\top, y^\top) * \Sigma^{-1} * \xy - \frac{1}{2} y^\top * \Sigma_Y^{-1} * y}
		\end{align*}
		Mit \cref{it: ue1-aufgabe3a} erhalten wir
		\begin{equation*}
			(x^\top, y^\top) * \Sigma^{-1} * \xy - y^\top * \Sigma_Y^{-1} * y
			= \schlange{x} (\Sigma / \Sigma_Y)^{-1} \schlange{x}
		\end{equation*}	
		mit $\schlange{x} = x - \Sigma_{XY}\Sigma_Y^{-1} y$. Definieren wir $\schlange{\mu} \defeq \Sigma_{XY}\Sigma_Y y$, so hast $f_{X|Y=y}$ die Form
		\begin{equation*}
			f_{X|Y=y}(x) = c * \exp\brackets{-\frac{1}{2} (x - \schlange{\mu}) * (\Sigma / \Sigma_Y)^{-1} * (x - \schlange{\mu})}
		\end{equation*}
		Dies ist die Dichte einer Realisation von $X|Y$. Für $X|Y$ erhalten wir folglich
		\begin{equation*}
			\EW[X | Y=y] = \schlange{\mu} \follows \mu = \EW[X | Y] = \Sigma_{XY} \Sigma_Y^{-1} 
			\quad \und \quad
			\Cov(X | Y) = (\Sigma / \Sigma_Y)
		\end{equation*}
	\end{enumerate}
	
	\undef\xy
	
% AUFGABE 4
	\begin{task}
		Seien $(X_n)_{n \in \N}$ und $(Y_n)_{n \in \N}$ zwei Martingale bezüglich $(\F_n)_{n \in \N}$ mit $X_n, Y_n \in L^2$ für alle $n \in \N$. Beweisen Sie:
		\begin{enumerate}[label=(\alph*)]
			\item $\EW[X_m Y_n | \F_m] = X_m Y_m$ fast sicher für alle $m \le n$
			\item $\EW[X_n Y_n] - \EW[X_0 Y_0] = \sum_{j=1}^n \EW[(X_j - X_{j-1})(Y_j - Y_{j-1})]$
			\item $\Var(X_n) = \Var(X_0) + \sum_{j=1}^n \Var(X-j - X_{j-1})$
			\item die Zufallsvariablen $X_0, X_1 - X_0, X_2 - X_1, \dots X_j - X_{j-1}$ sind paarweise orthogonal
		\end{enumerate}
	\end{task}
	
	\begin{enumerate}[label=(\alph*), leftmargin=*]
		\item \label{it: ue1-aufgabe4a} Seien $(X_n)_n$ und $(Y_n)_n$ Martingale bezüglich $F_n$. Da $X_m$ $\F_m$-messbar ist, gilt für alle $m \le n$
		\begin{equation*}
			\E[X_m Y_n | \F_m ] 
			= X_m * \E[Y_n\mid \F_m] 
			= X_m Y_m \quad \text{fast sicher}
		\end{equation*}
		\item \label{it: ue1-aufgabe4b} Für $1 \le j \le n$ betrachten wir
		\begin{align*}
			&\E[(X_j - X_{j-1})(Y_j - Y_{j-1})] \\
			= \enskip &\E[\E[X_j Y_j - X_j Y_{j-1} - X_{j-1} Y_j + X_{j-1} Y_{j-1} | \F_{j-1}]] \tag{Turmregel} \\
			= \enskip &\E[X_j Y_j] - \E[X_{j-1} Y_{j-1}] - \E[X_{j-1} Y_{j-1}] + \E[X_{j-1} Y_{j-1}] \tag{\ref{it: ue1-aufgabe4a} \& Turmregel} \\
			= \enskip &\E[X_j Y_j] - \E[X_{j-1} Y_{j-1}]
		\end{align*}
		Damit gilt
		\begin{equation*}
			\sum_{j=1}^n \E[(X_j - X_{j-1}) (Y_j - Y_{j-1})] 
			= \sum_{j=1}^n \E[X_j Y_j] - \E[X_{j-1} Y_{j-1}] 
			= \E[X_n Y_n] - \E[X_0 Y_0]
		\end{equation*} 
		mithilfe einer Teleskopsumme.
		\item Für Martingale $(X_n)_n$ gilt $\E[X_n] = \text{const.}$ für alle $n \in \N$, da $\EW[X_0] = \EW[{\EW[X_n | \F_0]}] = \EW[X_n]$ für alle $n \ge 0$. Damit gilt dann unter Nutzung des Verschiebunssatzes
		\begin{equation*}
			\Var(X_n) - \Var(X_0) = \E[X_n^2] - \E[X_n]^2 - \E[X_0^2] + \E[X_0]^2 = \E[X_n^2] - \E[X_0^2] 
		\end{equation*}
		sowie auch $\Var(X_j - X_{j-1}) = \E[(X_j - X_{j-1} -\E[X_j - X_{j-1}])^2] = \E[(X_j-X_{j-1})^2]$ für $1 \le j \le n$.
		Mit der Überlegung aus \cref{it: ue1-aufgabe4b} und $Y_j = X_j$ erhalten wir $\Var(X_j - X_{j-1}) = \E[X_j^2] - \E[X_{j-1}^2]$.
		Mit einem Teleskopsummenargument gilt dann
		\begin{equation*}
			\sum_{j=1}^n \Var(X_j - X_{j-1}) = \sum_{j=1}^n \E[X_j^2] - \E[X_{j-1}^2] = \E[X_n^2]-\E[X_0^2]
		\end{equation*}
		Damit ist also
		\begin{equation*}
			\Var(X_n) - \Var(X_0) = \sum_{j=1}^n \Var(X_j - X_{j-1})
		\end{equation*}
		woraus die Behauptung folgt.
		\item Seien $i,j \in \N$, wobei wir annehmen, dass $i \leq j$. Es gilt
		\begin{equation*}
			\scal{X_0}{X_j - X_{j-1}} 
			= \EW[X_0(X_j - X_{j-1})] 
			= \EW[X_0 X_j - X_0 X_{j-1}] 
			= \EW[X_0 X_j] - \EW[X_0 X_{j-1}]
		\end{equation*}
		Mit der Turmregel und \cref{it: ue1-aufgabe4a} erhalten wir
		\begin{equation*}
			\EW[X_0 X_j] - \EW[X_0 X_{j-1}]
			= \E[\EW[X_0 X_j | \F_0]] - \E[\EW[X_0 X_{j-1} | \F_0]]
			= \EW[X_0^2] - \EW[X_0^2] = 0
		\end{equation*}
		für alle $j \ge 1$.
		Außerdem gilt mit analogem Vorgehen
		\begin{align*}
			\scal{X_i- X_{i-1}}{X_j -X_{j-1}} 
			&= \E[(X_i- X_{i-1}) (X_j -X_{j-1})] \\
			&= \E[ X_i X_j - X_i X_{j-1} - X_{i-1} X_j + X_{i-1} X_{j-1}] \\
			&= \E[X_i^2] -\E[X_i^2]-\E[X_{i-1}^2]+\E[X_{i-1}^2] \tag{\ref{it: ue1-aufgabe4a} \& Turmregel} \\
			&= 0
		\end{align*}
	\end{enumerate}

\end{exercisePage}



\end{document}
